\documentclass[a4paper,11pt]{ltjsarticle}
\usepackage[haranoaji]{luatexja-preset}
\usepackage{amsmath,amssymb}
\usepackage[top=12mm,bottom=10mm,left=14mm,right=14mm]{geometry}
\usepackage{array}
\usepackage{tikz}
\usetikzlibrary{calc}
\pagestyle{empty}
\setlength{\parindent}{0pt}
\newlength{\qcol}\setlength{\qcol}{0.47\textwidth}
\newlength{\qgap}
\newlength{\anscolw}\setlength{\anscolw}{0.30\textwidth}
\newlength{\ansh}
\newcommand{\Q}[2]{\parbox[t]{\qcol}{\raggedright (#1)\hspace{0.6em}$\displaystyle #2$}}
\newcommand{\QR}[4]{\Q{#1}{#2}\hfill\Q{#3}{#4}\par\vspace{\qgap}}
\newcommand{\anscell}[1]{\rule[-\ansdrop]{0pt}{\ansh}(#1)}
\newlength{\ansdrop}

\begin{document}
{\large\bfseries 計算問題　24}\hfill 名前(\underline{\hspace{32mm}})\par\vspace{3mm}
■（1）〜（16）の計算をしなさい。（17）〜（21）は方程式を解きなさい。\par\vspace{3mm}
\setlength{\qgap}{5.4mm}
\QR{1}{-9^2+(-9)^2-4}{2}{9\div9\times(-36)}
\QR{3}{(-81)\div(+9)}{4}{-0.9-(-4.9)-1.8}
\QR{5}{(9x-9)-(8x+4)}{6}{\dfrac{9}{9}(9x-72)-\dfrac{4}{8}(8x+72)}
\QR{7}{\dfrac{9}{9}\div4+\dfrac{8}{9}}{8}{\dfrac{x-9}{9}\times36}
\QR{9}{\dfrac{1}{9}(81x-36)+\dfrac{1}{8}(32x+72)}{10}{(81a-36)\div9}
\QR{11}{(-16)+(-12)-(-9)}{12}{17-(-9)\div\left(-\dfrac{9}{4}\right)}
\QR{13}{(9a+9)+(4a-8)}{14}{(-13)\times(-11)}
\QR{15}{10\times\{9-(-4+1)\times2\}}{16}{9(9x-2)-(4x-8)}
\QR{17}{9x-9=-27}{18}{9(x+9)=10x+81}
\QR{19}{0.9x+0.9=0.0}{20}{\dfrac{5}{6}x-1=\dfrac{1}{2}x+2}
\vspace{1mm}
\Q{21}{\dfrac{9x-9}{4}-\dfrac{x+8}{4}=-1}\par\vspace{3mm}
\setlength{\ansh}{9.5mm}\setlength{\ansdrop}{6mm}%
\begin{center}\begin{tabular}{|p{\anscolw}|p{\anscolw}|p{\anscolw}|}\hline
\anscell{1} & \anscell{2} & \anscell{3} \\\hline
\anscell{4} & \anscell{5} & \anscell{6} \\\hline
\anscell{7} & \anscell{8} & \anscell{9} \\\hline
\anscell{10} & \anscell{11} & \anscell{12} \\\hline
\anscell{13} & \anscell{14} & \anscell{15} \\\hline
\anscell{16} & \anscell{17} & \anscell{18} \\\hline
\anscell{19} & \anscell{20} & \anscell{21} \\\hline
\end{tabular}\end{center}

\end{document}
